\documentclass[10pt,a4paper]{amsart}
\usepackage[margin=19mm]{geometry}
\usepackage{amsmath,amssymb,mathtools}
\usepackage[scr=rsfs,cal=euler]{mathalfa}
\usepackage{xfrac}
\usepackage[unicode,hidelinks,bookmarks=false]{hyperref}
\newcommand{\ie}{i.e.\ }
\newcommand{\cf}{cf.\ }
\newcommand{\resp}{respectively\ }
\newcommand{\Subsection}{\subsection}
\providecommand{\nobreakdash}{\nobreak\hbox{-}}
\setlength{\emergencystretch}{2em}
\multlinegap0pt
\usepackage{bm}
\usepackage{bbm}
\usepackage{dsfont}
\usepackage{xspace}
\usepackage{cancel}
\usepackage{mathtools}
\usepackage{amsthm}
\makeatletter
\@ifundefined{thm}{\newtheorem{thm}{Thm}}{}
\@ifundefined{example}{\newtheorem{example}[thm]{Example}}{}
\@ifundefined{lem}{\newtheorem{lem}[thm]{Lemma}}{}
\@ifundefined{rmk}{\newtheorem{rmk}[thm]{Remark}}{}
\makeatother
\makeatletter
\expandafter\ifx\csname customThm\endcsname\relax
\newenvironment{customThm}[1]{\paragraph*{\textbf{Theorem #1.}}\itshape}{\par}
\fi
\newcommand*{\defeq}{\mathrel{\vcenter{\baselineskip0.5ex \lineskiplimit0pt
			\hbox{\scriptsize.}\hbox{\scriptsize.}}}%
	=}
\makeatother
\title[Analyticity of the Hausdorff dimension and metric structures]{Analyticity of the Hausdorff dimension and metric structures on Misiurewicz families of polynomials}
\author{Fabrizio Bianchi, Yan Mary He}
\date{Source: arXiv:2508.12493v1 (CC BY 4.0)}
\begin{document}
\maketitle

\noindent\emph{Source and licence.} Analyticity of the Hausdorff dimension and metric structures on Misiurewicz families of polynomials. Version arXiv:2508.12493v1 (CC BY 4.0), distributed under the Creative Commons Attribution 4.0 International licence, as recorded in the arXiv licence metadata for this exact version and shown on the abstract page https://arxiv.org/abs/2508.12493v1. Full rights notice: licenses/arxiv-2508.12493-CC-BY-4.0.txt.\par\smallskip

\noindent\textit{Editorial scope.} Counted unit: the Lem whose source label is \texttt{l:L4MS} in the article (source0.tex lines 1159--1171, source label \texttt{l:L4MS}), delivered together with the article's own complete original proof of it (source0.tex lines 1173--1371, one \texttt{proof} environment of 199 lines). No mathematics is written, repaired or re-derived: statement and proof are the article's own text line for line. Required context retained uncounted: ex:deg3 (lines 644--668); ex:deg3-preim (lines 752--768); l:L2MS (lines 813--819); l:Hlambda-logq (lines 845--851). Every definition, display and label of those lines is retained unchanged. Omitted from this package and not redistributed: 1--643, 669--751, 769--812, 820--844, 852--1158, 1172--1172, 1372--2435. Nothing inside a retained excerpt is omitted or altered: every retained body is delivered exactly as the article's lines are written, so no omission receipt of any kind accompanies this package. Citations: the retained excerpts carry no citation command at all, so no bibliography entry of the article is transcribed and no reference is redistributed. Reference adaptation: none was needed and none was made; every reference inside the delivered unit resolves inside it, so no unresolved reference or citation is produced. Printed slips kept literal: no printed word, symbol, hypothesis or number of the counted statement or of its proof was altered. Layout and class: the article's own class is \texttt{amsart} with packages including inputenc, biblatex, amssymb, amsfonts, amsmath, amscd, xy, enumerate, mathrsfs, graphicx, fontenc, color, hyperref, dsfont; the wrapper is the lane's pinned \texttt{amsart} editorial wrapper at 10pt on A4, which restates the theorem-like environments the retained text needs and adds the article's own macro definitions that the retained text needs (\texttt{\textbackslash defeq}), with extra wrapper packages (bm, bbm, dsfont, xspace, cancel, mathtools). The delivered numbering is the unit's own. Assets: no retained span contains a figure environment, a graphics-inclusion command, a PDF-page inclusion, a table or a TikZ picture; no image file is copied into this package, the package declares no asset dependency and no image file is redistributed with it. Licence: version arXiv:2508.12493v1 of this article is distributed under the Creative Commons Attribution 4.0 International licence, as recorded in the arXiv licence metadata for this exact version and shown on the abstract page https://arxiv.org/abs/2508.12493v1. A full rights notice is delivered at licenses/arxiv-2508.12493-CC-BY-4.0.txt. Characters: every retained excerpt is pure ASCII, so the delivered engine, XeLaTeX with its default Latin Modern text font, reports no missing character.
\begin{example}\label{ex:deg3}
Assume that the degree of $F$ is 
equal to $3$. In this case, $P_1$ is given by $U_1\cup V_1$,
where
$U_1 = P_1 (c)$
and
$V_1=P_1(r)$. The set $P_2$ is given by the following five sets: 
\begin{itemize}
    \item  $U_2=P_2(c)$, which is mapped by $F$ to $V_1$ with degree $2$, 
    \item $V_2 \defeq P_2 (r)$,
    which is mapped univalently by $F$ to $V_1$;
    \item a topological disc $W_2\subset V_1\setminus V_2$, which is mapped univalently by $F$ to $U_1$, and
    \item two topological disks $Y_{2,1}, Y_{2,2}
    \subset U_1\setminus U_2$, each mapped univalently
    by $F$ to $U_1$.
\end{itemize}
More generally,
for every $n\geq 2$, there exist
two components $Y_{n,1}, Y_{n,2}\subset D_{n-1}$
and a component $W_n\subset V_{n-1}\setminus V_{n}$ of $P_n$
which are all
mapped 
univalently
by $F^n$ to $U_1$. 
\end{example}

\begin{example}\label{ex:deg3-preim}
Consider again the case where $d=3$ as in Example \ref{ex:deg3}. We see that every point not on the first floor of $\mathcal T$ 
has exactly one preimage under $T$ 
(given by the corresponding point on the floor just below it). On the other hand, on the first floor:
\begin{itemize}
\item every point in $V_1$ has exactly one preimage in $V_2$;
\item every point in $U_1$ has infinitely many preimages: one in $W_2$, and then, for every $n\geq 2$, one in each $Y_{n,1}$ and $Y_{n,2}$.
\end{itemize}

\begin{rmk}\label{rmk_bdryUk}
Observe that all the points of the form $(z,k)$ with $z$ on the 
boundary
of $U_k$ are mapped
(to the first floor and) to the boundary of $V_1$. In particular, their images are not in $\mathcal T$.
\end{rmk}

\end{example}

\begin{lem}\label{l:L2MS}
There is a constant $C$ independent of $n$ and $\lambda$ such that if $T_\lambda^n$
is defined and differentiable at $y$, then
\[
\frac{|(T^n_\lambda)'' (y)|}{|(T^n_\lambda)' (y)|^2} \leq C.
\]
\end{lem}

\begin{lem}\label{l:Hlambda-logq}
Up to taking $N(0)$ and $\kappa$ sufficiently small, 
the family $\{H_\lambda\}_{\lambda \in\overline{ N(0)}}$
is compact
in the space 
of $\kappa$-H\"older continuous maps.
\end{lem}

\begin{lem}\label{l:L4MS}
For all 
$t_1,t_2 \in \mathbb R$ with $t_1+t_2>0$
and $\lambda_1,\lambda_2\in N(0)$,
there exist a sequence $c_n$ with $0 < c_n\to 0$ 
and a constant $C>0$ such that for
every $n\ge1$, and every $g \in C^{\kappa}(\mathcal T)$, 
we have
$$\|\mathcal L^n_{t_1,\lambda_1, t_2, \lambda_2}
g\|'_{\kappa} \le c_n
\eta(t_1,\lambda_1, t_2, \lambda_2)^n
\|g\|'_{\kappa} +C \|\mathcal L^n_{t_1,\lambda_1, t_2, \lambda_2}|g| \|_{\infty}.$$
\end{lem}

\begin{proof}
As above, for ease of notation, we will only give the proof for the case where $d=3$ and we will simply
denote $\mathcal L_{t_1, \lambda_1, t_2, \lambda_2}$
by $\mathcal L$.
We also only consider the restriction to $U_1 \times \{1\}$, as the estimates are simpler otherwise (since only one preimage is involved and the system is expanding).

\medskip

The $n$-th iterate of $\mathcal L$ is given by
\[
\mathcal{L}
^n g(x) = \sum_{T^n(y) = x} g(y) \cdot
\left(
\prod_{j=0}^{n-1}
|R_{\lambda_1}(T^{j}(y))|^{-t_1}
%\big|
|R_{\lambda_2}(T^{j}(y))|^{-t_2}
\right)
\quad \mbox{ for all }
g \in C(\mathcal T).
\]
Hence, 
we have
\[
\begin{aligned}
\|\mathcal L^n
g\|'_{\kappa} &=
\sup_{\substack{a,b\in \mathcal T}}
\frac{1}{d_{\mathcal{T}}(a,b)^{\kappa}} \cdot|\mathcal L^n
g(a)-
\mathcal L^n
g(b)|\\
& =\sup_{\substack{a,b\in \mathcal T}}
\frac{1}{d_{\mathcal{T}}(a,b)^{\kappa}} \cdot
\Big|\sum_{T^n(y_a) = a}  g(y_a)
\Big(
\prod_{j=0}^{n-1}
|R_{\lambda_1}(T^{j}(y_a))|^{-t_1}
|R_{\lambda_2}(T^{j}(y_a))|^{-t_2}
\Big)
\\
& \quad  \quad -\sum_{T^n(y_b) = b} g(y_b)
%|R_\lambda(y_b)\cdots R_\lambda(T^{n-1}y_b)|^{-t}
\Big(
\prod_{j=0}^{n-1}
|R_{\lambda_1}(T^{j}(y_b))|^{-t_1}
%\big|
|R_{\lambda_2}(T^{j}(y_b))|^{-t_2}
\Big)
\Big|.
\end{aligned}\]


We now recall (see Example \ref{ex:deg3-preim}) that we can pair the $n$-preimages of $a$ and $b$ under $F^n$
depending on the set of $P_n$ to which they belong. We will then denote these preimages by $\{a_i\}$ and $\{b_i\}$,
meaning that, for every $i_0$, $a_{i_0}$
and $b_{i_0}$ belong to the same component of $P_n$. The above expression is then bounded by $I+ II+III+IV$, where



\begin{align*}
I & \defeq \sup_{\substack{a,b\in \mathcal T}}
\frac{1}{d_{\mathcal{T}}(a,b)^{\kappa}} \cdot \sum_{i} |g(a_i)-g(b_i)|
\cdot
\Big(
\prod_{j=0}^{n-1}
|R_{\lambda_1}(T^{j}(a_i))|^{-t_1}
|R_{\lambda_2}(T^{j}(a_i))|^{-t_2}
\Big),\\
II & \defeq
\sup_{\substack{a,b\in \mathcal T}}
\frac{1}{d_{\mathcal{T}}(a,b)^{\kappa}} \cdot \sum_{i}|g(b_i)|\cdot  \Big|
\prod_{j=0}^{n-1}
|R_{\lambda_1}(T^{j}(b_i))|^{-t_1}
-
\prod_{j=0}^{n-1}
|R_{\lambda_1}(T^{j}(a_i))|^{-t_1}
\Big|
\cdot
\prod_{j=0}^{n-1}
|R_{\lambda_2}(T^{j}(b_i))|^{-t_2},
\\
III & \defeq
\sup_{\substack{a,b\in \mathcal T}}
\frac{1}{d_{\mathcal{T}}(a,b)^{\kappa}} \cdot \sum_{i}|g(b_i)|\cdot 
\prod_{j=0}^{n-1}
|R_{\lambda_1}(T^{j}(b_i))|^{-t_1}
\cdot
\Big|
\prod_{j=0}^{n-1}
|R_{\lambda_2}(T^{j}(b_i))|^{-t_2}
-
\prod_{j=0}^{n-1}
|R_{\lambda_2}(T^{j}(a_i))|^{-t_2}
\Big|,\quad \mbox{ and }\\
IV & \defeq
\sup_{\substack{a,b\in \mathcal T}}
\frac{1}{d_{\mathcal{T}}(a,b)^{\kappa}} \cdot \sum_{i}|g(b_i)|\cdot 
\prod_{j=0}^{n-1}
|R_{\lambda_1}(T^{j}(b_i))|^{-t_1}
-
\prod_{j=0}^{n-1}
|R_{\lambda_1}(T^{j}(a_i))|^{-t_1}
\Big|
\cdot\\
& \quad \quad
\quad \quad
\quad \quad
\quad \quad
\quad \quad
\quad \quad
\quad \quad
\quad \quad  \Big|
\prod_{j=0}^{n-1}
|R_{\lambda_2}(T^{j}(b_i))|^{-t_2}
-
\prod_{j=0}^{n-1}
|R_{\lambda_2}(T^{j}(a_i))|^{-t_2}
\Big|.
\end{align*}

For $I$, we have
\begin{align*}
I &
\le \|\mathcal{L}
^n 1\|_\infty \cdot \sup_{\substack{a,b\in \mathcal T}}
\frac{\max_{i} |g(a_i)-g(b_i)|}{d_{\mathcal{T}}(a,b)^{\kappa}}
\le \|\mathcal{L}
^n 1\|_\infty \cdot \sup_{\substack{a,b\in \mathcal T}}
\frac{\max_i \|g\|_{\kappa}' \cdot d_{\mathcal{T}}(a_i,b_i)^\kappa}{d_{\mathcal{T}}(a,b)^{\kappa}}  \\
& \le
\|\mathcal{L}
^n 1\|_\infty \cdot \|g\|_{\kappa}'\cdot c_n\lesssim
\eta(t_1,\lambda_1, t_2, \lambda_2)^n
\cdot\|g\|_{\kappa}'\cdot  c_n
\end{align*}
for some $c_n$ as in the statement.
In the third inequality, we have used the fact that $\max_i d_\mathcal{T}(a_i,b_i) \le d_{\mathcal{T}}(a,b)\cdot c'_n$ for some $c'_n>0$ 
with $c'_n \to 0$ as $n \to \infty$,
since $T$ is expanding.

\medskip

For $II$, observe that we have $II=0$ if $t_1=0$. Otherwise, we compute
\begin{align*}
II & =\sup_{\substack{a,b\in \mathcal T}}
\frac{1}{d_{\mathcal{T}}(a,b)^{\kappa}}
\cdot \sum_{i}|g(b_i)|\cdot  \left|
|(T_{\lambda_1}^n)'(H_{\lambda_1}(a_i))|^{-t_1} - |(T_{\lambda_1}^n)'(H_{\lambda_1}(b_i))|^{-t_1}
\right|
\cdot
\left|
(T^n_{\lambda_2})' (H_{\lambda_2} (b_i))
\right|^{-t_2}
\\
& \lesssim
 |t_1|
\cdot 
\sup_{\substack{a,b\in \mathcal T}} \frac{1}{d_{\mathcal{T}}(a,b)^{\kappa}} \cdot
\sum_{i}
|g(b_i)|
 \left|
\frac{
(T_{\lambda_1}^n)''(H_{\lambda_1}(b_i))}{(T_{\lambda_1}^n)'(H_{\lambda_1}(b_i))|^{t_1+2}}
 \right|
 \cdot
\left| (T_{\lambda_2}^n)'(H_{\lambda_2}(b_i))\right|^{-t_2}
 \cdot d_{\mathcal{T}}(H_\lambda(y_a),H_\lambda(y_b))
\\
 & \lesssim 
  |t_1|
  \cdot
\|\mathcal{L}^{n}
|g|\|_\infty \cdot\sup_{\substack{a,b\in \mathcal T}} \frac{1}{d_{\mathcal{T}}(a,b)^{\kappa}} \cdot \max_{i} \frac{|(T^n_\lambda)''(H_\lambda(b_i))|}{|(T^n_\lambda)'(H_\lambda(b_i))|^2}
 \cdot 
  d_{\mathcal{T}}(H_\lambda(a_i),H_\lambda(b_i))\\ 
 & \lesssim
|t_1| \cdot \|\mathcal{L}^{n}
 |g|\|_\infty \cdot\sup_{\substack{a,b\in \mathcal T}} \frac{1}{d_{\mathcal{T}}(a,b)^{\kappa}}
 \cdot
 d_{\mathcal{T}}(a,b)^\kappa
 = |t_1| \cdot \|\mathcal{L}^{n}
 |g|\|_\infty,
\end{align*}
where in the last 
line
we used 
Lemmas \ref{l:L2MS} and
\ref{l:Hlambda-logq}.
For $III$, we can use the same steps as for $II$,
up to reversing the role of $(t_1, \lambda_1)$
and $(t_2,\lambda_2)$. The term $IV$ is similarly
bounded by 
$\|\mathcal{L}^{n} |g|\|_\infty$
using the same estimates as in $II$
and Lemma \ref{l:L2MS}
on both the differences of products inside the sum (in particular, we have $IV=0$ if either $t_1=0$ or $t_2=0$).
The proof is complete. 
\end{proof}

\end{document}
